\section{Component, Prefab, dan Scene}

\begin{learningoutcomes}
\item Menambahkan dan mengkonfigurasi komponen Unity
\item Membuat dan menggunakan Prefab
\item Mengelola Scene dan Build Settings
\end{learningoutcomes}

\subsection{Menambahkan Komponen}

Rigidbody 2D untuk fisika (Gravity Scale = 0 untuk space). Box Collider 2D atau Polygon Collider 2D untuk deteksi tabrakan. Add Component di Inspector.

\subsection{Prefab}

\begin{definisi}[Prefab]
Prefab adalah template GameObject yang dapat digunakan kembali dan diinstansiasi berkali-kali \cite{UnityManual2024}.
\end{definisi}

Membuat Prefab: Buat GameObject (Bullet, Asteroid), tambahkan komponen, drag ke folder Prefabs, hapus dari Hierarchy. Gunakan untuk objek yang di-spawn berulang.

\subsection{Scene Management}

Simpan scene (Ctrl+S) di folder Scenes. File $\rightarrow$ New Scene untuk scene baru. File $\rightarrow$ Build Settings: tarik scene ke Scenes In Build, atur urutan (MainMenu index 0).
